\makeatletter
\def\input@path{{./}{../}{../../}{../../../}{../../../../}{preamble/}{../preamble/}{../../preamble/}{../../../preamble/}{../../../../preamble/}{book2026-2027/preamble/}}
\makeatother
\input{preamble_exams.tex}
\input{graphs/G11_Histograms.tex}
\graphicspath{{./}{../}{../../}{../../../}{../../../../}{preamble/}{../preamble/}{../../preamble/}{../../../preamble/}{../../../../preamble/}{book2026-2027/preamble/}}
\setYColumnWidth{0.27}
\setlength{\columnsep}{6mm}
\usepackage{needspace}
\NewCommandCopy\SharedSubsectionBox\SubsectionBox
% Preserve the shared numbered/starred forms within the local width adjustment.
\RenewDocumentCommand{\SubsectionBox}{s m}{%
  \begingroup\setlength{\linewidth}{\dimexpr\linewidth-2\fboxsep\relax}%
  \IfBooleanTF{#1}{\SharedSubsectionBox*{#2}}{\SharedSubsectionBox{#2}}\endgroup}
\hypersetup{pdftitle={Grade 11 Term 1 C1--C7 Catch-up Exam Version 2}}

\setlength{\fboxsep}{1pt}

\begin{document}
\noindent
\begin{tabularx}{\textwidth}{|C{2.8cm}|C{\dimexpr\textwidth-5.6cm-6\tabcolsep-4\arrayrulewidth\relax}|C{2.8cm}|}
\hline
\includegraphics[width=2.5cm]{logo.png} &
\CellCenter{\textbf{GLOBAL ECONOMICS}\par
\textbf{Grade 11}\par
\textbf{Catch-up exam -- V2}\par
\textbf{Continuous distributions and sample data}\par
\textbf{Teacher: Nicolás López Cuéllar}} &
\CellCenter{\textbf{1st TERM}\par\textbf{2026--2027}}\\\hline
\end{tabularx}

\LearningObjective
{Identify probability distributions and their supports; use geometric interval probabilities and exponential waiting times; calculate theoretical moments, empirical probabilities, and sample statistics; and interpret errors relative to theoretical models in economic and operational contexts.}
{	\footnotesize
	\textbf{C1.} Characterises discrete and continuous random variables, identifies suitable distributions and supports, and interprets density and cumulative probability in context.\par
	\textbf{C2.} Applies uniform and triangular models and constant or piecewise-linear densities to interval probabilities using geometric areas.\par
	\textbf{C3.} Uses circular geometry and arc proportions to calculate probabilities for uniform timing and directions around a circle.\par
	\textbf{C4.} Applies exponential models to constant-rate waiting times, infers event rates, and assesses model suitability.\par
	\textbf{C5.} Calculates and interprets the mean, variance, and standard deviation of continuous distributions using standard formulas.\par
	\textbf{C6.} Estimates interval probabilities using sample relative frequencies and compares them with theoretical probabilities.\par
	\textbf{C7.} Calculates sample means, variances, and standard deviations for different sample sizes and data representations and compares them with theory.
}

\StudentInfo \vspace{-3mm}

\begin{multicols}{2}

\SubsectionBox*{Criteria assessment}
This exam assesses C1--C7. Each question and task shows its points. \textbf{Passing requires a minimum of 5 cumulative points per criterion, including points earned in previous learning activities}. Show all procedures; answers without the required procedure receive no credit. Partial credit may be awarded for a correct method with arithmetic errors. Missing names or late submissions incur a 1-point deduction per criterion.

\begin{center}
	\small\setlength{\tabcolsep}{5pt}
	\begin{tabular}{|c|c|c|c|c||c|c|c|c|}
		\hline
		\textbf{\textnumero}&\textbf{C1}&\textbf{C2}&\textbf{C3}&\textbf{C4}&
		\textbf{\textnumero}&\textbf{C5}&\textbf{C6}&\textbf{C7}\\\hline
		1&6&--&--&--&8&2&--&--\\\hline
		2&--&6&--&--&9&2&--&--\\\hline
		3&--&--&3&--&10&2&--&--\\\hline
		4&--&--&3&--&11&--&2&--\\\hline
		5&--&--&--&3&12&--&2&--\\\hline
		6&--&--&--&2&13&--&2&--\\\hline
		7&--&--&--&2&14&--&--&2\\\hline
		&&&&&15&--&--&2\\\hline
		&&&&&16&--&--&2\\\hline
		\textbf{Tot.}&\textbf{6}&\textbf{6}&\textbf{6}&\textbf{7}&
		\textbf{Tot.}&\textbf{6}&\textbf{6}&\textbf{6}\\\hline
	\end{tabular}
\end{center}

\columnbreak

\raggedcolumns
\SubsectionBox{Models in Colombian operations \score{\mbox{C1-\sml{6}}}}
\textbf{For each variable below}, classify it as discrete or continuous, identify its distribution, parameters and support, and briefly justify your choice.

\textbf{1. Municipal compost participation.} A programme records (a) whether a randomly selected household sorts organic waste correctly (1 if correct, 0 otherwise); and (b) the number sorting correctly among 35 households. Households sort correctly independently with probability 0.86. \score{1}

\textbf{2. Water-meter auditing.} A utility records (a) the number of independent tests up to and including the first faulty meter, with constant fault probability 0.06; and (b) a calibration error on $[-2,4]$ litres, with equal probabilities for equal-length intervals and no errors outside this range. Give (a) the probability of the first fault on test $k$ and (b) the density function. \score{1}

\textbf{3. Cultural-centre bookings.} A centre records (a) the number of rehearsal bookings in one hour: arrivals are independent at a constant rate of 3 per hour, one at a time, with negligible multiple arrivals in sufficiently short intervals; and (b) room-reset time on $[8,26]$ minutes: its density rises linearly from zero at 8 to a single peak at 15, then falls linearly to zero at 26, and is zero elsewhere. \score{1}

\textbf{4. Solar-energy monitoring.} A cooperative records (a) the usable battery-bank fraction at sunset on $[0.2,0.9]$: its density rises linearly from zero at 0.2 to its maximum at 0.9 and is zero elsewhere; and (b) sensor error with an exact symmetric bell curve on the real line, centred at 1 watt with standard deviation 3 watts. For (a), find the normalised density function and maximum height. \score{1}

\textbf{5. Agricultural training.} A service records (a) the number of cancellations among 45 independently booked places, each cancelled with probability 0.07; and (b) the number of registration calls over four days: arrivals are independent at a constant rate of 6 per day, one at a time, with negligible multiple arrivals in sufficiently short intervals. \score{1}

\textbf{6. Community-kitchen planning.} A kitchen studies (a) vegetable-batch preparation time from 12 to 36 minutes, with short times most common and density decreasing steadily to zero at the upper limit; and (b) daily rice use, symmetric around 18 kg, with values farther from 18 less common and standard deviation 2.5 kg. For each, discuss whether the distribution and support are uniquely determined. Propose a model, state any additional assumptions, and distinguish model support from possible physical values. Give the density function for (a) under your assumption. \score{1}


\SubsectionBox{Seedling-tray preparation \score{\mbox{C2-\sml{6}}}}
At a Boyac\'a nursery, tray-filling and seeding time $S$ has support $[4,28]$ minutes. Its density rises linearly from zero at 4 to a single maximum at 14, then falls linearly to zero at 28; it is zero elsewhere.

Identify the model and support, sketch the density, find its maximum height and full piecewise formula (including outside the support), and calculate $P(8\le S\le22)$ using geometric areas. \score{3}

Change the support to $[2,34]$ minutes, keeping zero density at both endpoints and linear sides meeting at the mode of 14 minutes. Identify the new model, sketch its density, find its maximum height and full piecewise formula, and recalculate $P(8\le S\le22)$ using geometric areas. \score{3}



\SubsectionBox{A repeating irrigation cycle \score{\mbox{C3-\sml{3}}}}
A flower farm's pressure fluctuation occurs uniformly within a repeating 48-minute irrigation cycle, starting at minute 0 and restarting immediately after minute 48.

Find the probability of a fluctuation from minute 41 through the restart to minute 9 of the next cycle. State the event as interval(s) within one cycle. \score{1}

Find the probability of a fluctuation within 150 seconds on either side of the cycle start. State the event as interval(s) within one cycle. \score{2}


\begin{minipage}{\linewidth}

\SubsectionBox{Direction of a floating current marker \score{\mbox{C3-\sml{3}}}}
A river-monitoring team's floating marker has an initial direction uniform around a full circle. A bank marker indicates north. Angular tolerances apply on \textbf{both sides} of north, including wraparound.

For each tolerance ($\pi/8$ radians and $18^\circ$ from north), state the event as angular intervals with north at zero and calculate its probability. Identify the larger probability, give the difference in percentage points, and calculate the relative increase from the $18^\circ$ probability. \score{3}
\end{minipage}\par


\SubsectionBox{Public-library return arrivals \score{\mbox{C4-\sml{3}}}}
Returns enter a public library's automated chute as a Poisson process, with a different constant rate in each scenario below. Let $T$ be the waiting time in minutes from a fixed observation time to the next return.

For each scenario, identify the waiting-time model, rate in returns per minute and support, and calculate the probability. Give the general density function and briefly justify the model.

With an average of 7 returns every 35 minutes, find $P(T<3)$. \score{0.5}

With a mean waiting time of 5 minutes between returns, find $P(T>8)$. \score{0.5}

With an average of 4 returns every 25 minutes, find $P(1<T<6)$. \score{1}

With an average of 18 returns per hour, find $P(T<2\text{ or }T>7)$. \score{1}


\SubsectionBox{Bicycle-hire arrival rates \score{\mbox{C4-\sml{2}}}}
At a municipal bicycle-hire kiosk, rentals start as a Poisson process with different constant morning and afternoon rates. Let $T_m,T_a$ be the respective waiting times in minutes from a fixed time in each shift to the next rental. Assume each shift continues throughout the wait.

In the morning, $P(T_m>12)=0.35$. Infer the rate in rental starts per minute and per hour. \score{1}

In the afternoon, $P(T_a<9)=0.72$. Infer the rate in rental starts per minute and per hour. \score{1}


\SubsectionBox{Two recycling collection gates \score{\mbox{C4-\sml{2}}}}
Vehicle entries at recycling gates A and B follow independent constant-rate Poisson processes. From a fixed observation time, A's next entry has probability 0.84 of occurring within 4 minutes; B's has probability 0.60 of taking more than 4 minutes. Count both gates' entries from this same time.

Infer both rates in entries per minute and state the expected-count functions. Find when A\textquotesingle s expected count exceeds B\textquotesingle s by 4, in minutes, and explain what this time means for actual entries. \score{2}




\SubsectionBox{Two document-digitisation desks \score{\mbox{C5-\sml{2}}}}
At a municipal archive, $A$ and $B$ are file-digitisation times in minutes at two desks, both with support $[2,17]$. A's density rises linearly from zero at 2 to its maximum at 17. B has equal probabilities for equal-length intervals within the support. Both densities are zero elsewhere.

Identify A\textquotesingle s distribution, parameters and support, and calculate its mean, variance and standard deviation. Interpret the mean in context. \score{1}

Identify B\textquotesingle s distribution, parameters and support, and calculate its mean, variance and standard deviation. Interpret the mean and standard deviation in context. \score{1}


\SubsectionBox{\small Pump maintenance and daily operating costs \score{\mbox{C5-\sml{2}}}}
A rural water association's pump maintenance duration $D$ has support $[2,14]$ hours. Its density falls linearly from its maximum at 2 to zero at 14 and is zero elsewhere.

Daily operating cost $R$, in \textbf{thousands of Colombian pesos (COP)}, has an exact symmetric bell-shaped density on the real line, centred at 2400 with inflection points 2250 and 2550.

Identify D\textquotesingle s distribution, parameters and support, and calculate its mean, variance and standard deviation. Interpret the mean in context. \score{1}

Identify R\textquotesingle s distribution and infer its parameters and support. Calculate mean and SD in \textbf{COP} and variance in \textbf{COP$^2$}. Interpret the mean. \score{1}



\SubsectionBox{\small Museum admissions and audio-guide sessions \score{\mbox{C5-\sml{2}}}}
Museum admissions follow a constant-rate Poisson process averaging 15 per hour. Let $T$ be the interarrival time in minutes.

Audio-guide session duration $S$ has support $[6,27]$ minutes. Its density rises linearly from zero at 6 to a single peak at 15, then falls linearly to zero at 27; it is zero elsewhere.

Identify T\textquotesingle s distribution, parameters and support, and calculate its mean and SD in minutes and its variance in minutes$^2$. Interpret the mean in context. \score{1}

Identify S\textquotesingle s distribution, parameters and support, and calculate its mean and SD in minutes and its variance in minutes$^2$. Interpret the mean and standard deviation in context. \score{1}


\SubsectionBox{Commercial-kitchen water use \score{\mbox{C6-\sml{2}}}}
A catering kitchen records water use $X$ in litres for 50 cleaning batches in two non-overlapping parts. The first 20 observations are:
\[\begin{array}{rrrrr}4.1&4.5&5.0&5.3&5.8\\6.0&6.2&6.7&7.0&7.4\\4.8&5.0&5.5&5.9&6.4\\6.8&7.0&7.2&7.6&8.0\end{array}\]
The remaining 30 observations have this discrete frequency table:
\[\begin{array}{c|rrrrrr}x_i&5.0&5.4&5.8&6.3&7.0&7.5\\\hline f_i&4&6&5&7&3&5\end{array}\]
Combine both parts as one sample. The model gives $P(5\le X<7)=0.58$.

State the sample size and event count, then calculate the empirical probability. \score{1}

Calculate its absolute and percentage errors. \score{1}


\columnbreak

\SubsectionBox{Produce cooling, grouped data \score{\mbox{C6-\sml{2}}}}
Produce-crate cooling times $X$ in minutes at a packing centre are grouped below. Classes include the left endpoint and exclude the right.
\[\begin{array}{c|rrrr}\text{Class}&[2,7)&[7,12)&[12,17)&[17,22)\\\hline f_i&4&7&6&4\end{array}\]
The model gives $P(7\le X<17)=0.68$.

State the sample size and event count, then calculate the empirical probability. \score{1}

Calculate its absolute and percentage errors. \score{1}


\SubsectionBox{Battery-swap service, histogram \score{\mbox{C6-\sml{2}}}}
The histogram is the only sample record of battery-swap times $X$ at an electric-motorcycle station. The horizontal axis shows minutes; bar heights are frequencies. Equal-width classes include the left endpoint and exclude the right.
\HistogramGraph{(1,4),(4,7),(7,10),(10,13)}{3,9,7,6}{0.94\linewidth}{4.4cm}{blue!25}{3}{1}
The model gives $P(4\le X<10)=0.60$.

State the sample size and event count, then calculate the empirical probability. \score{1}

Calculate its absolute and percentage errors. \score{1}



\SubsectionBox{Daily paper use, mixed data \score{\mbox{C7-\sml{2}}}}
A legal-aid office records paper use in hundreds of sheets on 13 days in two non-overlapping parts: the first five observations are $13,16,18,21,23$; the remaining eight are summarised below:
\[\begin{array}{c|rrr}x_i&13&18&23\\\hline f_i&1&4&3\end{array}\]
Combine both parts as one sample. The theoretical mean is 18 hundreds of sheets and variance is 16 (hundreds of sheets)$^2$.

State the sample size; calculate the sample mean, variance and standard deviation. Show your frequency summary and weighted sum of squared deviations. \score{1}

Calculate absolute and percentage errors for mean and variance. \score{1}




\SubsectionBox{\small Compost maturation checks, grouped data \score{\mbox{C7-\sml{2}}}}
An urban garden groups times between compost checks in days. Classes include the left endpoint and exclude the right. The theoretical mean is 15 days and variance is 24 days$^2$.
\[\begin{array}{c|rrr}\text{Class}&[5,11)&[11,17)&[17,23)\\\hline f_i&3&2&2\end{array}\]

State the sample size; calculate the sample mean, variance and standard deviation. Show your frequency summary and weighted sum of squared deviations. \score{1}

Calculate absolute and percentage errors for mean and variance. \score{1}


\SubsectionBox{Community-bakery flour use, histogram \score{\mbox{C7-\sml{2}}}}
The histogram is the only sample record of flour use per batch at a community bakery. The horizontal axis shows kilograms; bar heights are frequencies. Equal-width classes include the left endpoint and exclude the right. The theoretical mean is 8 kg and variance is 7 kg$^2$.
\HistogramGraph{(3,6),(6,9),(9,12),(12,15)}{3,2,3,2}{0.94\linewidth}{4.4cm}{orange!30}{3}{1}

State the sample size; calculate the sample mean, variance and standard deviation. Show your frequency summary and weighted sum of squared deviations. \score{1}

Calculate absolute and percentage errors for mean and variance. \score{1}


\end{multicols}
\end{document}
